% !TeX root = ../main.tex
\section{Исследование анализатора спектра в режиме обзора}
\subsection{Фиксация начальных параметров прибора (Preset)}
Перед началом проведения измерений на входной разъем анализатора спектра RIGOL DSA1030 был установлен защитный адаптер с разделительным конденсатором и аттенюатор на $20\text{ дБ}$. Прибор был включен и переведен в исходное состояние (Preset).

Начальные параметры развертки по частоте, полосы фильтрации и уровень собственных шумов, отобразившиеся на экране анализатора спектра, занесены в табл.~\ref{tab:initial_settings}.

\begin{table}[H]
    \caption{Начальные установки параметров анализатора спектра RIGOL DSA1030}
    \label{tab:initial_settings}
    \begin{tabular}{l c c}
        \toprule
        \textbf{Параметр} & \textbf{Обозначение} & \textbf{Значение} \\
        \midrule
        Центральная частота & \textit{Center Freq} & $1{,}5\text{ ГГц}$ \\
        Полоса обзора & \textit{Span} & $3{,}0\text{ ГГц}$ \\
        Полоса пропускания фильтра ПЧ & \textit{RBW} & $1\text{ МГц}$ (Auto) \\
        Полоса видеофильтра & \textit{VBW} & $1\text{ МГц}$ (Auto) \\
        Опорный уровень & \textit{Ref Level} & $0\text{ дБм}$ \\
        Входное ослабление & \textit{Input Atten} & $10\text{ дБ}$ \\
        Отображаемый средний уровень шума & \textit{DANL} & $-32\text{ дБм}$ \\
        \bottomrule
    \end{tabular}
\end{table}
По полученным начальным установкам видно, что прибор по умолчанию переходит в обзорный режим, охватывающий весь диапазон рабочих частот (от $9\text{ кГц}$ до $3\text{ ГГц}$). 

Широкая полоса пропускания фильтра промежуточной частоты ($\text{RBW} = 1\text{ МГц}$) и включённый входной аттенюатор обеспечивают высокую скорость первичного сканирования тракта. При этом отображаемый средний уровень шума ($\text{DANL}$) составляет $-32\text{ дБм}$, что говорит об испроавности коаксиального кабеля. Интегральная мощность тепловых шумов прямо пропорциональна ширине полосы пропускания тракта ПЧ:
\begin{equation}
    P_{\text{ш}} = G_0 \cdot \text{RBW},
\end{equation}
где $G_0$ --- спектральная плотность мощности шума. 

Исходный вид экрана анализатора спектра с установленной шумовой дорожкой и параметрами развёртки приведён на рис.~\ref{fig:fig_1}.
\begin{figure}[H]
    \centering
    \includegraphics[width=0.9\linewidth]{fgi1}
    \caption{Стартовые значения параметров на экране анализатора спектра}
    \label{fig:fig_1}
\end{figure}

\subsection{Параметры анализатора спектра в режиме AutoTune}

На генераторе был задан режим непрерывных гармонических колебаний с выбранной рабочей частотой несущей $f_0 = 1{,}0\text{ ГГц}$ и выходной мощностью $0\text{ дБм}$.

После подачи сигнала на панели управления анализатора была нажата клавиша \textbf{AutoTune}. Прибор автоматически определил частоту входного сигнала, отцентрировал её на экране, скорректировал полосу обзора, полосы фильтрации и опорный уровень. Полученные значения параметров приведены в табл.~\ref{tab:autotune_settings}.


\begin{table}[H]
    \caption{Параметры установки в режиме AutoTune}
    \label{tab:autotune_settings}
    \begin{tabular}{l c c}
        \toprule
        \textbf{Параметр} & \textbf{Обозначение} & \textbf{Значение} \\
        \midrule
        Центральная частота & \textit{Center Freq} & $999{,}98\text{ МГц}$ \\
        Полоса обзора & \textit{Span} & $10{,}0\text{ ГГц}$ \\
        Полоса пропускания фильтра ПЧ & \textit{RBW} & $300\text{ кГц}$ \\
        Полоса видеофильтра & \textit{VBW} & $300\text{ кГц}$ \\
        \bottomrule
    \end{tabular}
\end{table}
